Ceahl\u{a}u Mountain is located in Neam\c{t} County, not far from the city of Piatra Neam\c{t}. An observer on the top of the mountain ($\varphi=46^\circ53'57''$ N, $L=26^\circ24'41''$ E, altitude $h=1907\,\mathrm m$ above sea level), looks at the starry sky, making observations of the star Almach in the constellation Andromeda.
\begin{description}
\item[(a)] Determine whether, from the top of Mount Ceahl\u{a}u, the star Almach ($\gamma$ AND, $\alpha=02^{\mathrm h}03^{\mathrm m}54^{\mathrm s}$, $\delta=42^\circ19'47''$) is circumpolar. The radius of the Earth is known to be $R_\oplus=6378\,\mathrm{km}$.
\item[(b)] In 1778, it was discovered that $\gamma$ AND is a double star. The two stars are separated by an angle of about $10''$. What should be the diameter of the objective lens of a telescope to see the two components distinctly (the wavelength is known to be $\lambda=550\,\mathrm{nm}$)?
\item[(c)] It was later discovered that the Almach star is a multiple star system consisting of the primary component $\gamma_1$ (A) with apparent magnitude $m_1=2.27$ and a subsystem $\gamma_2$ with apparent magnitude $m_2=4.84$. Calculate the apparent magnitude of the system consisting of the star $\gamma_1$ and the subsystem $\gamma_2$.
\item[(d)] What is the absolute magnitude of the star, knowing that it is approximately 355 light-years away?
\end{description}
